\documentclass[10pt,journal,compsoc]{IEEEtran}

\usepackage{graphicx}
\usepackage{amsmath}
\usepackage{url}
\usepackage{booktabs}
\usepackage{listings}
\usepackage{xcolor}

\begin{document}

\title{Demand-Driven Context: A Framework for Building Enterprise AI Agent Knowledge Bases Through Problem-Driven Curation}

\author{
  \IEEEauthorblockN{[Author Name]}
  \IEEEauthorblockA{[Affiliation]\\
  [Email]}
}

\markboth{IEEE Software}%
{Demand-Driven Context}

\IEEEtitleabstractindextext{%
\begin{abstract}
Enterprise AI agents excel at reasoning but lack the domain-specific context needed to make accurate decisions about real systems, integrations, and architectures. Current approaches---context window stuffing, retrieval-augmented generation, and top-down knowledge curation---fail because they attempt to anticipate what knowledge an agent needs rather than letting actual problems reveal it.

We present Demand-Driven Context (DDC), a framework inspired by Test-Driven Development where real enterprise problems drive knowledge curation. In DDC, an AI agent attempts to solve a problem with zero domain context, identifies its information gaps via a structured checklist, receives targeted answers from domain experts, and curates the validated knowledge into a structured, navigable repository. We apply DDC in [an enterprise domain] across [N] problems over [M] months, measuring knowledge base growth, context reuse across problems, agent accuracy improvement, and curation time per cycle.

Our results show [coverage convergence / entity reuse / time reduction] after [N] cycles, suggesting that 20-30 real problems can build a more focused and maintainable knowledge base than months of top-down documentation. We also report emergent patterns including self-organizing learning paths, a curator agent pattern, and unprompted agent identity discovery.
\end{abstract}

\begin{IEEEkeywords}
AI agents, knowledge management, context engineering, demand-driven development, enterprise software
\end{IEEEkeywords}}

\maketitle

\IEEEdisplaynontitleabstractindextext

\section{Introduction}
% The context problem for enterprise AI agents
% Code generation vs. cognitive work
% Why new-joiner onboarding is the right analogy
% Contribution summary

\textit{[To be written after evidence collection --- Phase 5 of roadmap]}

\section{Background and Related Work}
% RAG and its limitations for enterprise context
% Active learning and failure-driven knowledge acquisition
% Knowledge engineering (traditional and modern)
% Context engineering for LLMs
% Lean/pull systems as a conceptual ancestor
% TDD as structural parallel

\textit{[To be written after literature review --- Gate 1]}

\section{The DDC Method}
% Formalized loop: problem -> sandbox -> checklist -> human answers -> curation -> re-attempt -> graduation
% The information checklist artifact
% Knowledge base structure (entities, relationships, frontmatter)
% The convergence hypothesis
% TDD structural mapping

\textit{[To be written --- primary contribution section]}

\section{Evaluation}
% Setting: enterprise domain, N teams, M-year transformation
% Metrics: coverage curve, curation time, agent accuracy, context reuse
% Data collection: structured cycle logs
% Results
% Comparison: no context vs. top-down vs. DDC

\textit{[To be written after 20+ cycles --- Phase 2 of roadmap]}

\section{Discussion}
% When DDC works and when it doesn't
% The curator agent pattern
% Emergent learning paths
% Agent self-discovery
% Threats to validity
% Scalability considerations

\textit{[To be written]}

\section{Conclusion and Future Work}
% Summary of contributions
% Future: multi-domain evaluation, federated context, curator agent formalization

\textit{[To be written]}

\bibliographystyle{IEEEtran}
\bibliography{references}

\end{document}